\tikzstyle{obj1}=[line width = 0.8pt, line join=round, line cap=round]
\tikzstyle{obj2}=[line width = 0.8pt, line join=round, line cap=round, color=BrickRed]
\tikzstyle{obj3}=[line width = 0.8pt, line join=round, line cap=round, color=BlueViolet]
\tikzstyle{surf}=[line width = 0.8pt, line join=round, line cap=round, fill=black, fill opacity=.06, rounded corners=2pt]
\tikzstyle{mpstensor}=[draw=black, line width = 0.7pt, fill=BlueViolet!14, circle, inner sep=3pt]
\tikzstyle{anyontensor}=[draw=black, line width = 0.7pt, fill=orange!40, circle, inner sep=3pt]
\tikzstyle{mpotensor}=[draw=black, line width=.7pt, fill=gray!20, rectangle, rounded corners=2pt, inner sep=3.8pt]
\tikzstyle{matrix}=[draw=black, line width = 0.7pt, fill=black!10, circle, inner sep=1.6pt]
\tikzstyle{mor}=[draw=none, circle, fill=gray!22]
\tikzstyle{fusiontensor} = [draw=black, line width=.7pt, fill=gray!20, circle, inner sep=2.9pt]
\tikzset{mpoint/.style={ transform shape, sloped, 
		minimum width=20pt, minimum height=4pt, inner sep=0pt, ultra thin, font=\tiny}}
\tikzstyle{dot}=[dash pattern= on 0.6pt off 2.1pt]

\tikzset{wiggly/.style={decorate, decoration={snake, segment length=4.1pt, amplitude=1.pt}}}

\def\centerarc[#1](#2)(#3:#4:#5)% Syntax: [draw options] (center) (initial angle:final angle:radius)
{ \draw[#1] ($(#2)+({#5*cos(#3)},{#5*sin(#3)})$) arc (#3:#4:#5); }


\tikzset{->1-/.style={decoration={
			markings,
			mark=at position #1 with {\arrow{Triangle[fill=black, width=3.5pt, length=3.5pt]}}},postaction={decorate}}
		}
\tikzset{-<1-/.style={decoration={
			markings,
			mark=at position #1 with {\arrowreversed{Triangle[fill=black, width=3.5pt, length=3.5pt]}}},postaction={decorate}}
		}
\tikzset{->2-/.style={decoration={
			markings,
			mark=at position #1 with {\arrow{Triangle[fill=BrickRed, width=3.5pt, length=3.5pt]}}},postaction={decorate}}
		}
\tikzset{-<2-/.style={decoration={
			markings,
			mark=at position #1 with {\arrowreversed{Triangle[fill=BrickRed, width=3.5pt, length=3.5pt]}}},postaction={decorate}}
		}
\tikzset{->t-/.style={decoration={
			markings,
			mark=at position #1 with {\arrow{Triangle[fill=gray!40, width=3.5pt, length=3.5pt]}}},postaction={decorate}}
		}
\tikzset{->tt-/.style={decoration={
			markings,
			mark=at position #1 with {\arrow{Triangle[fill=gray!40, width=3.5pt, length=3.5pt]}},mark=at position {#1+.05} with {\arrow{Triangle[fill=gray!40, width=3.5pt, length=3.5pt]}}},postaction={decorate}}
}
%\tikzset{->3-/.style={decoration={
%			markings,
%			mark=at position #1 with {\arrow{Triangle[fill=BlueViolet!80, width=3.5pt, length=3.5pt]}}},postaction={decorate}}}
%
\tikzset{
	mask point/.style={ transform shape, sloped, 
		minimum width=20pt, minimum height=4pt, inner sep=0pt, ultra thin, font=\tiny}
}
\tikzset{
	clip even odd rule/.code={\pgfseteorule},
	invclip/.style={clip,insert path=[clip even odd rule]{
			[reset cm](-\maxdimen,-\maxdimen)rectangle(\maxdimen,\maxdimen)
}}}
\newcommand\clipIntersection[1]{
	(#1.north west) -- (#1.north east) -- (#1.south east) -- (#1.south west) -- (#1.north west)
}
\newcommand\clipAroundNode[2]{
	\begin{pgfscope}
		\begin{scope}[overlay]
			\path[invclip] \clipIntersection{#1};#2
		\end{scope}
	\end{pgfscope}
}
\newcommand\clipAroundTwoNodes[3]{
	\begin{pgfscope}
		\begin{scope}[overlay]
			\path[invclip] \clipIntersection{#1} \clipIntersection{#2};#3
		\end{scope}
	\end{pgfscope}
}
\newcommand\clipAroundThreeNodes[4]{
	\begin{pgfscope}
		\begin{scope}[overlay]
			\path[invclip] \clipIntersection{#1} \clipIntersection{#2} \clipIntersection{#3};#4
		\end{scope}
	\end{pgfscope}
}
\newcommand\clipAroundFourNodes[5]{
	\begin{pgfscope}
		\begin{scope}[overlay]
			\path[invclip] \clipIntersection{#1} \clipIntersection{#2} \clipIntersection{#3} \clipIntersection{#4};#5
		\end{scope}
	\end{pgfscope}
}
\newcommand{\middleCoo}[6]{
	\path (#1) --coordinate[pos=0.5](#3) (#2);
	\path (#1) --coordinate[pos={0.5-\sp}](#4) (#2);
	\path (#1) --coordinate[pos={0.5+\sp}](#5) (#2);
}
\newcommand{\barycentre}[8]{
	\middleCoo{#1}{#2}{A}{}{}{}{#8};
	\middleCoo{#2}{#3}{B}{}{}{}{#8};
	\middleCoo{#1}{#3}{C}{}{}{}{#8};
	\path[name path=PA#3] (A) -- (#3);
	\path[name path=PB#1] (B) -- (#1);
	\path[name path=PC#2] (C) -- (#2);
	\path [name intersections={of=PA#3 and PB#1,by={#4}}];
	\path (#1) --coordinate[pos=0.84](#5) (#4);
	\path (#2) --coordinate[pos=0.84](#6) (#4);
	\path (#3) --coordinate[pos=0.84](#7) (#4);
}
\newcommand{\barycentreSq}[9]{
	\middleCoo{#1}{#2}{A1}{A2}{A3}{};
	\middleCoo{#3}{#4}{B1}{B2}{B3}{};
	\middleCoo{#1}{#3}{C1}{C2}{C3}{};
	\middleCoo{#2}{#4}{D1}{D2}{D3}{};
	\middleCoo{A1}{B1}{#5}{}{};
	\path[name path=PA2B2] (A2) -- (B2);
	\path[name path=PA3B3] (A3) -- (B3);
	\path[name path=PC2D2] (C2) -- (D2);
	\path[name path=PC3D3] (C3) -- (D3);       
	\path [name intersections={of=PA2B2 and PC2D2,by={#6}}];
	\path [name intersections={of=PA2B2 and PC3D3,by={#8}}];
	\path [name intersections={of=PA3B3 and PC2D2,by={#7}}];
	\path [name intersections={of=PA3B3 and PC3D3,by={#9}}];
}